\chapter{Deques}\label{ch:deques}

\section{The Deque ADT}

\begin{definition}[Deque]
A \emph{deque} (double-ended queue, pronounced ``deck'') is a sequence in which
insertions and deletions may take place at both ends. Its operations are
\cc{push\_front(x)}, \cc{push\_back(x)}, \cc{pop\_front()}, \cc{pop\_back()},
\cc{front()}, \cc{back()}, \cc{empty()} and \cc{size()}.
\end{definition}

\para{How we write it in C++} Exactly like \cc{deque<int>} from \cc{<deque>}:
\begin{lstlisting}[style=starter]
deque<int> d;
d.push_back(1);  d.push_front(2);     // d is 2 1
int a = d.front();  int b = d.back(); // a = 2, b = 1 (nothing removed)
d.pop_front();   d.pop_back();        // remove (they return nothing)
d.empty();  d.size();  d[i];          // d[i] reads item i
\end{lstlisting}

A deque generalises both earlier structures. Using only \cc{push\_back} and
\cc{pop\_back} it is a stack; using only \cc{push\_back} and \cc{pop\_front} it is a
queue. Two restricted forms appear in textbooks: an \emph{input-restricted} deque
(insertion at one end only, deletion at both) and an \emph{output-restricted} deque
(insertion at both ends, deletion at one only).

\begin{example}
Starting from an empty deque, the operations \cc{push\_back(1)}, \cc{push\_front(2)},
\cc{push\_back(3)}, \cc{push\_front(4)} produce \cc{4 2 1 3}. Then \cc{back()} is 3 and
\cc{front()} is 4; \cc{pop\_back()} and \cc{pop\_front()} remove them, leaving \cc{2 1}.
\end{example}

\section{The circular-array deque}

The circular buffer of Chapter~\ref{ch:queues} already supports \cc{push\_back} and
\cc{pop\_front} in $\Theta(1)$. The other two operations are the mirror images:
\begin{itemize}
\item \cc{push\_front(x)}: step \cc{front} \emph{back} by one cell, wrapping from 0 to
  \cc{CAP-1}, then write \cc{x} there and increase \cc{size};
\item \cc{pop\_back()}: the last item is at \cc{(front+size-1)\%CAP}; read it and
  decrease \cc{size} (\cc{front} does not change).
\end{itemize}

\listingfile[label=lst:deque]{lib/deque.h}{A circular-array deque written as a C++ class
(file \cc{deque.h}). Same functions as \cc{deque<int>}.}

\begin{note}{The modulo trap}
In C and C++ the result of \cc{\%} has the sign of the left operand, because integer
division truncates toward zero: \cc{-7 / 2 == -3} and \cc{-1 \% 5 == -1}, not $4$.
Hence \cc{front = (front - 1) \% CAP} sets \cc{front} to $-1$ when \cc{front} was $0$,
and the next write lands outside the array. Always add \cc{CAP} first:
\cc{(front - 1 + CAP) \% CAP}, or equivalently \cc{(front + CAP - 1) \% CAP}. With a fixed
capacity 5 the same step is often written \cc{(front + 4) \% 5}
(Tracing~\ref{trace:21}). In Python, \cc{-1 \% 5} is $4$; do not carry that habit
into C.
\end{note}

\begin{figure}[ht]\centering
\begin{tikzpicture}[x=7.5mm]
\foreach \v/\s [count=\i from 0] in {6/sh,{}/,{}/,{}/,{}/,{}/} \node[cell,\s] (a\i) at (\i,0) {\v};
\foreach \i in {0,...,5} \node[idx,above=0pt of a\i] {\i};
\draw[ptr] (0,-0.9) node[below,lbl]{front=0} -- (a0.south);
\node[left,font=\scriptsize] at (-0.6,0) {\texttt{[6]}};
\foreach \v/\s [count=\i from 0] in {6/sh,{}/,{}/,{}/,{}/,8/shd} \node[cell,\s] (b\i) at (\i+9,0) {\v};
\foreach \i in {0,...,5} \node[idx,above=0pt of b\i] {\i};
\draw[ptr] (14,-0.9) node[below,lbl]{front=(0-1+6)\%6=5} -- (b5.south);
\node[above,font=\scriptsize] at (11.5,0.6) {after \texttt{push\_front(8)}: \texttt{[8 6]}};
\end{tikzpicture}
\caption{\cc{push\_front} on a circular deque of capacity 6 whose front is cell 0: the
new front is the \emph{last} cell.}
\end{figure}

\section{The ArrayDeque: positional access in a circular array}

The course notes (after Morin's \emph{ArrayDeque}) go one step further and support the
whole List interface --- \cc{get(i)}, \cc{set(i,x)}, \cc{add(i,x)}, \cc{remove(i)} --- in
a circular array with fields \cc{j} (the index of item 0) and \cc{n} (the size). Item $i$
lives at \cc{a[(j+i)\%CAP]}.

\para{Shift the shorter side} To insert at position $i$, an array list moves the $n-i$
items to the right of $i$. In a circular array there is a second option: move the $i$
items to the \emph{left} of $i$ one cell to the left, after stepping \cc{j} back by one.
The ArrayDeque chooses the side with fewer items: the left side if $i<n/2$, otherwise the
right side. Removal is symmetric.

\listingfile[label=lst:arraydeque]{sol/p56.cpp}{ArrayDeque with positional insertion and
removal (Problem~\ref{prob:56}, Tracing~\ref{trace:22}).}

\begin{theorem}[Cost of ArrayDeque]
\cc{get} and \cc{set} cost $\Theta(1)$. \cc{add(i,x)} and \cc{remove(i)} move at most
$\min(i,\,n-i)$ items and cost $O(1+\min(i,n-i))$. In particular all four operations at
either end ($i=0$ or $i=n$, resp.\ $i=n-1$) are $\Theta(1)$, and the worst case, in the
middle, is about $n/2$ moves --- half that of an array list.
\end{theorem}
\begin{proof}
If $i<n/2$ the code moves the $i$ items before position $i$; otherwise $i\ge n/2$ and it
moves the $n-i$ items from position $i$ onward. In both cases the number moved is
$\min(i,n-i)$ up to one, and each move is one assignment.
\end{proof}

\section{The linked deque}

A singly linked list cannot be a deque with all operations in $\Theta(1)$: \cc{pop\_back}
needs the predecessor of the tail. A doubly linked list can. With a \emph{sentinel}
(Section~\ref{sec:sentinel}) closed into a circle --- \cc{sentinel->next} is the first
item, \cc{sentinel->prev} the last, and both point to the sentinel itself when the deque
is empty --- every insertion is ``link a node between two neighbours'' and every removal
is ``unlink a node from its two neighbours''. No pointer is ever \cc{NULL}, so the four
operations have no special cases at all.

\listingfile[label=lst:linkeddeque]{sol/p51.cpp}{A linked deque with a sentinel node
(Problem~\ref{prob:51}, Tracing~\ref{trace:25}).}

\begin{note}{Copying a linked deque copies only the pointer}
\cc{LinkedDeque e = d;} copies the two members: the \emph{pointer} \cc{sentinel} and the
number \cc{count}. Both objects now use the same nodes. A \cc{push} on \cc{d} is seen
when walking \cc{e}, but \cc{e.count} is not updated, so \cc{e.size()} becomes wrong
(Tracing~\ref{trace:25}). This is a \emph{shallow copy}. Do not copy a linked
structure unless you write code that copies every node.
\end{note}

\begin{table}[ht]\centering\small
\begin{tabular}{@{}lccc@{}}\toprule
 & circular array & ArrayDeque & doubly linked + sentinel\\\midrule
push/pop at either end & $\Theta(1)$ & $\Theta(1)$ & $\Theta(1)$\\
\cc{get(i)}, \cc{set(i,x)} & $\Theta(1)$ & $\Theta(1)$ & $\Theta(\min(i,n-i))$\\
\cc{add(i,x)}, \cc{remove(i)} & --- & $\Theta(1+\min(i,n-i))$ & $\Theta(\min(i,n-i))$ + $\Theta(1)$ link\\
memory per item & one cell & one cell & item + two pointers\\\bottomrule
\end{tabular}
\caption{Deque implementations. A doubly linked list walks from the nearer end to reach
position $i$.}
\end{table}

\section{The monotonic deque}\label{sec:monodeque}

\para{Problem} Given an array of $n$ values and a window size $k$, report the maximum of
every window \cc{a[i-k+1..i]}. Recomputing each maximum costs $\Theta(nk)$. A queue of
the window is not enough: when the maximum leaves the window, the next maximum is unknown.

\para{Idea} Keep a deque of \emph{indices} of the current window whose values
\emph{decrease} from front to back. When \cc{a[i]} arrives:
\begin{enumerate}
\item if the front index has slid out of the window ($\le i-k$), pop it from the
  \emph{front};
\item pop from the \emph{back} every index $j$ with $\cc{a[j]}\le\cc{a[i]}$: such $j$
  is older \emph{and} not larger than $i$, so it can never again be the maximum of a
  window that contains $i$;
\item push $i$ at the back; the front is now the index of the maximum of the window.
\end{enumerate}
The back behaves like the monotonic stack of Section~\ref{sec:monostack}, the front like
the queue of a sliding window. A deque is exactly the structure that does both.

\listingfile[label=lst:slidingmax]{sol/p53.cpp}{Sliding window maximum
(Problem~\ref{prob:53}, Tracing~\ref{trace:23}).}

\begin{table}[ht]\centering\small
\begin{tabular}{@{}ccllc@{}}\toprule
$i$ & \cc{a[i]} & popped & deque after (indices: values) & output\\\midrule
0 & 1 & --- & 0:1 & \\
1 & 3 & back 0 & 1:3 & \\
2 & $-1$ & --- & 1:3\; 2:$-1$ & 3\\
3 & $-3$ & --- & 1:3\; 2:$-1$\; 3:$-3$ & 3\\
4 & 5 & front 1; back 3, 2 & 4:5 & 5\\
5 & 3 & --- & 4:5\; 5:3 & 5\\
6 & 6 & back 5, 4 & 6:6 & 6\\
7 & 7 & back 6 & 7:7 & 7\\\bottomrule
\end{tabular}
\caption{Sliding window maximum on $[1,3,-1,-3,5,3,6,7]$ with $k=3$.}
\end{table}

\begin{theorem}[Sliding window maximum]
The algorithm is correct and runs in $\Theta(n)$ time and $O(k)$ extra space.
\end{theorem}
\begin{proof}
\emph{Invariant}: after round $i$ the deque contains, in increasing order of index,
exactly those indices $j$ of the window $[i-k+1,i]$ for which \cc{a[j]} is
\emph{greater} than every later value of the window; their values therefore decrease
from front to back, and the front holds the window maximum (the maximum is greater than
every later value, ties going to the latest). Step~1 removes the only index that can have
left the window, because the window advances by one. Step~2 removes exactly the indices
whose value is $\le\cc{a[i]}$, which are the ones no longer greater than every later
value. Step~3 adds $i$, trivially greater than every later value.
\emph{Cost}: each index is pushed once and popped at most once (from one end or the
other), so the loops of all rounds together do at most $2n$ deque operations. The deque
never holds more than $k$ indices.
\end{proof}

\para{Variants}
\begin{itemize}
\item \emph{Window minimum}: reverse the comparison; values increase from front to back.
\item \emph{Variable window} (two pointers): the window \cc{a[left..right]} grows at the
  right and shrinks at the left. Keeping one decreasing deque (maximum) and one
  increasing deque (minimum) answers ``longest subarray with $\max-\min\le
  \mathit{limit}$'' in $\Theta(n)$ (Problem~\ref{prob:54}).
\item \emph{Dynamic programming over a window}: when
  $\mathit{best}[i]=a[i]+\max(\mathit{best}[i-k..i-1])$, the deque keeps the indices of
  the last $k$ values of $\mathit{best}$, and the whole table costs $\Theta(n)$ instead
  of $\Theta(nk)$ (Jump Game VI, Problem~\ref{prob:57}).
\end{itemize}

\section{Further applications}

\para{Palindromes} Push the characters into a deque; while at least two are left, pop one
from each end and compare (Problem~\ref{prob:52}). The deque makes the ``two
pointers from both ends'' pattern explicit.

\para{Two deques back to back} A sequence split into a left half and a right half, each
in its own deque, allows $\Theta(1)$ access to the \emph{middle} as well as to the ends:
the middle is the back of the left deque or the front of the right one. After every
operation one item at most has to cross the boundary to restore the balance
$|\mathit{right}|-|\mathit{left}|\in\{0,1\}$ (front-middle-back queue,
Problem~\ref{prob:55}). The same trick with a single middle pointer on a linked list is
Problem~\ref{prob:24}.

\para{Undo with a limit} An editor that keeps only the last $100$ actions pushes new
actions at the back, undoes from the back (a stack), and discards the oldest action from
the \emph{front} when the limit is exceeded. Neither a stack nor a queue suffices; a
deque does.

\section{The standard library}\label{sec:stl}

In C++ the containers below implement the ADTs of this volume. Examinations that permit
the STL accept them; examinations that ask you to \emph{implement} an ADT do not.

\begin{table}[ht]\centering\small
\begin{tabular}{@{}lp{0.42\linewidth}p{0.31\linewidth}@{}}\toprule
Container & Main operations (all $\Theta(1)$ unless noted) & Notes\\\midrule
\cc{vector<int>} & \cc{push\_back}, \cc{pop\_back}, \cc{back}, \cc{[i]}, \cc{size};
  \cc{insert}/\cc{erase} at an iterator $\Theta(n)$ & amortised push; growth invalidates
  iterators\\
\cc{list<int>} & \cc{push\_front/back}, \cc{pop\_front/back}, \cc{front}, \cc{back},
  \cc{insert}/\cc{erase} at an iterator, \cc{splice} & doubly linked with sentinel; no
  \cc{[i]}\\
\cc{forward\_list<int>} & \cc{push\_front}, \cc{pop\_front}, \cc{insert\_after},
  \cc{erase\_after} & singly linked; no \cc{size()}\\
\cc{deque<int>} & \cc{push\_front/back}, \cc{pop\_front/back}, \cc{front}, \cc{back},
  \cc{[i]}; middle \cc{insert} $\Theta(n)$ & blocks of arrays\\
\cc{stack<int>} & \cc{push}, \cc{pop}, \cc{top}, \cc{empty}, \cc{size} & adapter over
  \cc{deque}\\
\cc{queue<int>} & \cc{push}, \cc{pop}, \cc{front}, \cc{back}, \cc{empty}, \cc{size} &
  adapter over \cc{deque}\\\bottomrule
\end{tabular}
\caption{STL containers.\label{tab:stl} In all of them \cc{pop}, \cc{pop\_front} and \cc{pop\_back}
return \cc{void}: read \cc{top()}/\cc{front()}/\cc{back()} first, and never call them on
an empty container (undefined behaviour, not an error code).}
\end{table}

\begin{table}[ht]\centering\small
\begin{tabular}{@{}p{0.42\linewidth}p{0.51\linewidth}@{}}\toprule
If the problem says \dots & think of \dots\\\midrule
both ends, front and back & a deque\\
maximum / minimum of every window & monotonic deque of indices (\ref{prob:53})\\
longest / shortest window with a max--min condition & two monotonic deques + two
pointers (\ref{prob:54})\\
$\mathit{best}[i]$ from the best of the last $k$ & DP + monotonic deque (\ref{prob:57})\\
compare from both ends & pop front and back (\ref{prob:52})\\
operations in the middle & two balanced deques (\ref{prob:55})\\
design a circular deque & the modulo trap (\ref{prob:50})\\\bottomrule
\end{tabular}
\caption{Recognising a deque problem.}
\end{table}
